\documentclass[10pt,a4paper]{article}
\usepackage[T1]{fontenc}
\usepackage[utf8]{inputenc}
\usepackage[ngerman]{babel}
\usepackage{lmodern}
\usepackage[margin=8mm]{geometry}
\usepackage{amsmath,amssymb,mathtools}
\usepackage{multicol}
\usepackage{array,booktabs,tabularx}
\usepackage{xcolor}
\usepackage{enumitem}
\usepackage{microtype}
\usepackage{fancyhdr}

\definecolor{Navy}{HTML}{17324D}
\definecolor{Blue}{HTML}{2F6FA3}
\definecolor{PaleBlue}{HTML}{EAF3FA}
\definecolor{Green}{HTML}{2E7D61}
\definecolor{PaleGreen}{HTML}{EAF6F1}
\definecolor{Orange}{HTML}{D9772A}
\definecolor{PaleOrange}{HTML}{FFF2E5}
\definecolor{Ink}{HTML}{1E2933}
\definecolor{Muted}{HTML}{5F6B76}

\setlength{\parindent}{0pt}
\setlength{\parskip}{1.2pt}
\setlength{\columnsep}{4mm}
\setlength{\columnseprule}{0.2pt}
\setlist[itemize]{leftmargin=3.2mm,itemsep=0pt,topsep=1pt}
\pagestyle{fancy}
\fancyhf{}
\renewcommand{\headrulewidth}{0pt}
\fancyfoot[C]{\scriptsize\color{Muted}Kompakte Formelsammlung \textbar{} Seite \thepage}

\newcommand{\CourseName}{Kursname}
\newcommand{\SheetName}{Kompakte Formelsammlung}

\newcommand{\titlebar}{%
  \noindent\colorbox{Navy}{%
    \parbox{\dimexpr\linewidth-2\fboxsep\relax}{%
      \color{white}\bfseries\large \CourseName\hfill\normalsize \SheetName}}
  \vspace{2mm}}

\newcommand{\sect}[1]{%
  \vspace{1.5mm}\noindent\colorbox{Blue}{%
    \parbox{\dimexpr\linewidth-2\fboxsep\relax}{\color{white}\bfseries #1}}%
  \vspace{0.8mm}}

\newcommand{\subsect}[1]{\vspace{0.8mm}\textbf{\color{Navy}#1}\par}

\newcommand{\directbox}[1]{%
  \noindent\colorbox{PaleGreen}{%
    \parbox{\dimexpr\linewidth-2\fboxsep\relax}{%
      \textbf{\color{Green}Direkt einsetzen.} #1}}\par}

\newcommand{\warnbox}[1]{%
  \noindent\colorbox{PaleOrange}{%
    \parbox{\dimexpr\linewidth-2\fboxsep\relax}{%
      \textbf{\color{Orange}Randfall.} #1}}\par}

\newenvironment{cheattable}
  {\renewcommand{\arraystretch}{1.05}\tabularx{\linewidth}{@{}>{\raggedright\arraybackslash\bfseries}p{14mm}>{\raggedright\arraybackslash}X@{}}}
  {\endtabularx}

\begin{document}
\titlebar
\begin{multicols*}{3}

\sect{Themenblock}
\subsect{Erkennung und Voraussetzungen}
Den Aufgabentyp kurz benennen. Voraussetzungen und lokale Symbolbedeutungen vor die Formel setzen.

\[
  y = f(x), \qquad x\in D
\]

\directbox{Gegebenes einsetzen, nach der gesuchten Größe umstellen und Einheit oder Dimension prüfen.}

\warnbox{Eine wichtige Ausnahme oder eine leicht verwechselte Voraussetzung knapp benennen.}

\subsect{Ablauf}
\begin{enumerate}[leftmargin=4mm,itemsep=0pt,topsep=1pt]
  \item Aufgabensignal erkennen.
  \item Voraussetzungen prüfen.
  \item Formel in Kursnotation anwenden.
  \item Ergebnis plausibilisieren.
\end{enumerate}

\sect{Vergleich}
\begin{cheattable}
Fall A & Kriterium, Formel und kurze Kontrolle.\\
Fall B & Abweichendes Kriterium und Randfall.\\
\end{cheattable}

\sect{Klausurcheck}
\begin{itemize}
  \item Sind alle Größen definiert?
  \item Stimmen Einheit, Vorzeichen und Wertebereich?
  \item Wurde der relevante Randfall geprüft?
\end{itemize}

\end{multicols*}
\end{document}
